% 来源：OEBPS/Text/01_46.xhtml（OPF spine 顺序；仅第一册）
\cleardoublepage
\BookEntry[荣·切尔诺《洛克菲勒传：全球首富的创富秘诀》]{荣·切尔诺《洛克菲勒传：全球首富的创富秘诀》}

\BookCover{images/书单30洛克菲勒传：全球首富的创富秘诀.jpg}{[美] 荣·切尔诺著| 王恩冕等译| 华东师范大学出版社2013年版}

\begin{BookSummaryBox}
\textbf{史上第一位亿万富翁、美国最著名企业王朝创建人洛克菲勒的一生}
\par
\end{BookSummaryBox}

\subsubsection*{这本书谈什么？}

洛克菲勒是有史以来第一位亿万富翁、美国最著名的企业王朝的创建人。洛克菲勒的父亲是犯重婚罪、行为放荡的假药贩子，母亲是循规蹈矩的清教徒。洛克菲勒出生在乡间，通过创建美国最有实力、最令人生畏的垄断性企业标准石油公司而发迹，成为世界首富。这家被文人称作“章鱼”的托拉斯企业所提炼和销售的石油，几乎占当时美国同类产品总产量的90\%。
\par

洛克菲勒的故事反映了美国历史上的一个转折期，它记录了南北战争之后小型企业发展成巨型公司的剧烈变化过程。这一变化不可逆转地改变了美国。本书用极其客观和优美的笔法，第一次向我们详尽、公允、深入地展现了这位实业巨子的全貌。
\par

\subsubsection*{芒格为什么谈起这本书？}

在伯克希尔哈撒韦公司1999财年的年会上，芒格说：
\par

\begin{quotation}
我是传记和自传的无可救药的爱好者，这一点我必须承认。即便知道这一点，我还是喜欢上了这本讲洛克菲勒的书。他的一生、他创办的企业，以及她的管理风格，都统统值得仔细研究。这本书还是有趣的垄断案例分析。我经常发现自己面临的情势与标准石油和微软迥异。
\par
\end{quotation}

\subsubsection*{《洛克菲勒传：全球首富的创富秘诀》作者荣·切尔诺如是说}

“作为一个肩负职责的小大人，他养成了一种终身不变的自我夸大的责任感。”
\par

“记住别人告诉你的东西很重要，比你自己早已知道的东西更重要。”
\par

“在很多情况下，最困难的竞争不是来自于强大、聪明、保守的对手，而是来自那些垂死挣扎的人，这些人对自己的成本浑然不知，却无论如何也要把工厂开下去，否则就得破产。”
\par

“洛克菲勒认为沉默代表着力量：只有内心虚弱的人才会随口乱讲，对着记者喋喋不休，而谨慎的商人则守口如瓶。‘成功来自多听少说’和‘只说不做的人就像是长满荒草的花园’是洛克菲勒最喜欢的两则箴言。”
\par

“我们许多人没能取得大的成就……不正是因为我们没有做到集中精力，不能在适当的时候排除外界干扰，把全部精力集中在要做的事情上吗？”
\par

“行善本身就是最好的回报，寻求别人感激的人将只会感到莫大的悲哀。”
\par

“沉默难以取胜处，畅所欲言者必胜。”
\par
